La presente guía documenta el software desarrollado para el prototipo educativo \textbf{ELCO-Dealer}, realizado en el marco de la asignatura \textbf{Electrónica de Consumo} del Grado en Ingeniería de Tecnologías y Servicios de Telecomunicación de la Universidad Politécnica de Madrid.

El objetivo del documento es explicar de forma ordenada la arquitectura del código, la lógica de control implementada en Arduino y la relación entre el software embebido, los sensores, los motores y la interfaz de usuario. Además, se incluye una descripción de la plataforma web asociada al proyecto, concebida como escaparate técnico y comercial del sistema.

Esta guía complementa al manual de usuario y montaje: mientras aquel se centra en el ensamblaje y operación física del prototipo, este documento describe cómo está organizado el programa, qué responsabilidades tiene cada módulo y qué parámetros deben revisarse durante la calibración.
